% !TEX program = xelatex
\documentclass[a4paper]{article}
\usepackage{amsthm}
\usepackage{amssymb}
\usepackage{bm}
\usepackage{mathtools}
\usepackage[x11names]{xcolor}
\usepackage{xparse}
\usepackage{fontspec}
\usepackage{unicode-math}
\setromanfont{DovesType-Regular.otf}
\setsansfont{Andika}
\setmathfont{Asana Math}[Scale=1]

% \usepackage{pstricks}
\usepackage{varwidth}
\usepackage{siunitx}
\usepackage{graphicx}
\usepackage[margin=1cm]{geometry}
\usepackage[most]{tcolorbox}
\usepackage{pgfplots}
\pgfplotsset{compat=newest}
\tcbuselibrary{skins,xparse,poster,breakable}
% \usetikzlibrary{fadings}
\usetikzlibrary{calc, plotmarks, shapes, shapes.geometric, positioning, angles, intersections, quotes, through, patterns, turtle, arrows.meta}
\usetikzlibrary{decorations.markings}
% \usepackage{etoolbox}
\usepackage{tkz-euclide}
% \usepackage{xlop}
% \newcommand\hole[2]{#1}

%%%%%%%%%%%%%%%%%%%%%%%%%%%%%%%%%%%%%%%%%%%%%%%%%%%%%%%%%
\newcommand\markangle[9]{% origin X Y radius radiusmark mark colour opacity
%  % fill red circle offset-from-centre
  \begin{scope}
    \path[clip] (#1) -- (#2) -- (#3);
    \fill[color=#7,fill opacity=#8,draw=black,name path global=pcircle]  % global declaration required otherwise pcircle is not seen by the `named intersections=' lines below.
    (#1) circle (#4);
  \end{scope}
  % middle calculation
  \path[name path=line one] (#1) -- (#2);
  \path[name path=line two] (#1) -- (#3);
  \path[%
  name intersections={of=line one and pcircle, by={inter one}},
  name intersections={of=line two and pcircle, by={inter two}}
  ] (inter one) -- (inter two) coordinate[pos=#9] (place);
  % put mark
  \node at ($(#1)!#5!(place)$) {\scriptsize{#6}};
}

% \newcommand{\condSoln}[2]{\ifcsdef{r@#1}{#2}{}}

% \newcommand\fadingtext[3][]{%
%    \begin{tikzfadingfrompicture}[name=fading letter]
%      \node[text=transparent!0,inner xsep=0pt,outer xsep=0pt,#1] {#3};
%    \end{tikzfadingfrompicture}%
%    \begin{tikzpicture}[baseline=(textnode.base)]
%      \node[inner sep=0pt,outer sep=0pt,#1](textnode){\phantom{#3}};
%      \shade[path fading=fading letter,#2,fit fading=false]
%      (textnode.south west) rectangle (textnode.north east);%
%    \end{tikzpicture}%
% }

\definecolor{JISpurple}{RGB}{89,72,122}
\definecolor{JISivory}{RGB}{241,234,221}
\definecolor{JIStaupe}{RGB}{183,156,154}
\definecolor{PaleGreen}{RGB}{240,255,240} % 'Honeydew'

\AddToHook{shipout/background}{%
    \put (0in,-\paperheight){\includegraphics[width=\paperwidth,height=\paperheight]{images/R10V5-R.png}}%
}

\newcommand\numberthis{\addtocounter{equation}{1}\tag{\theequation}}

\newtcolorbox{MyOuterBox}{%
  enhanced,
  breakable,
  frame style=JISpurple,
  colback=JISivory,
  colframe=JISpurple,
  title={\includegraphics[width=0.9cm,height=0.9cm]{images/JIS Final Logo FA-02.png}\raisebox{3mm}{\Large{Maths Challenge}\hspace{26em} \Large{\bfseries\sffamily 9}}},
}

\newtcolorbox{MyInnerBox}[2][]{enhanced,%empty,
coltitle=JISpurple,colback=white,
fonttitle=\bfseries\sffamily,
attach boxed title to top left={yshift=-1.5mm},
boxed title style={empty, size=small, top=1mm, bottom=0pt},
varwidth boxed title=0.5\linewidth,
frame code={
  \path (title.east|-frame.north) coordinate (aux);
\path[draw=JISpurple, line width=0.5mm, rounded corners,fill=white]
(frame.west) |- ([xshift=-2.5mm]title.north east) to[out=0, in=180] ([xshift=7.5mm]aux)-|(frame.east)|-(frame.south)-|cycle;
},
title={#2},#1}

\newtcolorbox{MyInnerSplitBox}[2][]{enhanced,%empty,
bicolor,sidebyside,sidebyside align=top seam,
righthand width=5.0cm,colbacklower=white,
coltitle=JISpurple,colback=white,
fonttitle=\bfseries\sffamily,
attach boxed title to top left={yshift=-1.5mm},
boxed title style={empty, size=small, top=1mm, bottom=0pt},
varwidth boxed title=0.5\linewidth,
frame code={
  \path (title.east|-frame.north) coordinate (aux);
\path[draw=JISpurple, line width=0.5mm, rounded corners,fill=white]
(frame.west) |- ([xshift=-2.5mm]title.north east) to[out=0, in=180] ([xshift=7.5mm]aux)-|(frame.east)|-(frame.south)-|cycle;
},
title={#2},#1}


\newtcolorbox{MySolutionBox}{%
  enhanced,
  breakable,
  frame style=JISpurple,
  colback=PaleGreen, colframe=green,
  title={\Large Solution},
  drop fuzzy shadow,
  halign=left,
}

%%%%%%%%%%%%%%%%%%%%%%%%%%%%%%%%%%%%%%%%%%%%%%%%%%
\newtoggle{SOLUTION}
%%% Uncomment the appropriate line below to show solutions %%%
% \toggletrue{SOLUTION}
\togglefalse{SOLUTION}
%%%%%%%%%%%%%%%%%%%%%%%%%%%%%%%%%%%%%%%%%%%%%%%%%


%%%%%%%%%%%%%%%%%%%%%%%%%%%%%%%%%%%%%%%%%%%%%%%%%%
%%%%%%            DOCUMENT BEGINS           %%%%%%
%%%%%%%%%%%%%%%%%%%%%%%%%%%%%%%%%%%%%%%%%%%%%%%%%%
\begin{document}


  \begin{MyOuterBox}
    \iftoggle{SOLUTION}{Here are the full, or partial solutions.
    }{
      Welcome to this week's Maths Challenge!\\
      Solutions must be explained in detail, responses with just the answer may be ignored.\\
      Drop your solution in the box in the staffroom by Tuesday.
    }
       \begin{MyInnerBox}{Year 9 and below}
      If
      \begin{equation*}
        ax = by = cz = 5\qquad\qquad\text{  and  }\qquad\qquad \frac{1}{x} + \frac{1}{y} + \frac{1}{z} = 4
      \end{equation*}
      What is the value of \(a + b + c\)?
      \iftoggle{SOLUTION}{%conditional output begin
      \begin{MySolutionBox}
        Notice that we need to find a relationship between \(a\), \(b\) and \(c\), so we could try eliminating \(x\), \(y\) and \(z\). Using the first set of equations, we can rewrite them as:\par
        \begin{alignat*}{4}
          ax &= 5 &\qquad\qquad by &= 5 &\qquad\qquad cz &= 5\\
          \frac{1}{x} &= \frac{a}{5} & \frac{1}{y} &= \frac{b}{5} & \frac{1}{z} &= \frac{c}{5}
        \end{alignat*}
        Now we can substitute these into the second equation of the question:
        \begin{align*}
          \frac{a}{5} + \frac{b}{5} + \frac{c}{5} &= 4\\
          5\left(\frac{a}{5} + \frac{b}{5} + \frac{c}{5}\right) &= 5\times 4\\
          a + b + c &= 20
        \end{align*}
      \end{MySolutionBox}
    }{}%conditional output end
    \end{MyInnerBox}


    \vspace{0.4cm}
          \begin{MyInnerBox}{Year 10 and above}
     Alnitak and Mintaka are \SI{300}{\kilo\m} apart. They travel towards each other in a straight line. Each travels at their own constant speed.\par
     If they both set off at 9:00 a.m. then they meet at noon.\par
     If Alnitak sets off at 6:00 a.m. and Mintaka sets off at 10:00 a.m. they still meet at noon.\par
     What speed does Alnitak travel at and what speed does Mintaka travel at?
      \iftoggle{SOLUTION}{%conditional output begin
      \begin{MySolutionBox}
        Since their speeds are constant we can use the \[Speed = \frac{distance}{time}\] formula.\par
        Let \(s_{a}\) be Alnitak's speed and let \(s_{m}\) be Mintaka's speed.\par
        We will let \(d_{1}\) be the distance from Alnitak's starting point to the place they meet when they both depart at 9 a.m., and \(d_{2}\) will be the distance from Alnitak's starting point to the place they meet when they depart at 6 a.m. and 10 a.m.\par
        In the first case, distance \(d_{1}\) is covered by Alnitak in \(\SI{3}{\hour}\). Then Mintaka travels \(300-d_{1}\) in 3 hours.\par
        In the second case, distance \(d_{2}\) is travelled by Alnitak in \(\SI{6}{\hour}\). Then Mintaka covers \(300-d_{2}\) in 2 hours.\par\medskip
      \begin{minipage}{0.4\textwidth}
        \underline{Case 1: Both depart at 9 a.m.}
        \begin{flalign*}
          s_{a} &= \frac{d_{1}}{3} &\\
          s_{m} &= \frac{300-d_{1}}{3} &
          \shortintertext{Eliminate \(d_{1}\) using these two equations.}
          3s_{a} &= 300-3s_{m} &\\
          3s_{a} + 3s_{m} &= 300 & \\
          s_{a} + s_{m} &= 100 &\\
        \end{flalign*}
      \end{minipage}%
      \hfill
      \begin{minipage}{0.4\textwidth}
        \underline{Case 2: 6 a.m. and 10 a.m. departures}
        \begin{flalign*}
          s_{a} &= \frac{d_{2}}{6} & \\
          s_{m} &= \frac{300-d_{2}}{2} &
          \shortintertext{Eliminate \(d_{2}\) using these two equations.}
          6s_{a} &= 300 - 2s_{m} & \\
          6s_{a} + 2s_{m} &= 300 & \\
          3s_{a} + s_{m} &= 150 & \\
        \end{flalign*}
      \end{minipage}\par
          Now we have a pair of simultaneous equations:\par
        \[
          \begin{rcases}
            s_{a} + s_{m} &= 100  \\[0em]
            3s_{a} + s_{m} &= 150
          \end{rcases}
        \]
        Solving, we find that\par
        \indent Alnitak's speed \(s_{a} = \SI{25}{\km\per\hour}\), and\par
        Mintaka's speed \(s_{m} = \SI{75}{\km\per\hour}\).
      \end{MySolutionBox}
    }{}%conditional output end
    \end{MyInnerBox}


  \end{MyOuterBox}

%%%%%%%%%%%%%%%%%%%%%%%%%%%%%%%%%%%%%%%%%%%%%%%%%%
\end{document}



